%density_waves.tex
%
%written by tools/paper_numbers.py; do not edit by hand

\newcommand{\dwM}{2}   %inputs.py satellite m
\newcommand{\dwMs}{2\times10^{-10}}   %inputs.py satellite mass_final
\newcommand{\dwAs}{1.590}   %inputs.py satellite r
\newcommand{\dwRL}{1.005}   %fast/satellite.py r_L (L2)
\newcommand{\dwAcrossPercent}{25}   %r_L's place across the ring, percent from the inner edge
\newcommand{\dwTauOrbits}{2100}   %tau_0.01 in orbits (L6)
\newcommand{\dwOrbits}{3300}   %the run's length in orbits
\newcommand{\dwRunFraction}{1.5}   %run length over tau_0.01
\newcommand{\dwXout}{0.015}   %the ring outer edge's x
\newcommand{\dwXfront}{0.014}   %the measured wave front's x at the run's end
\newcommand{\dwkZ}{0.01}   %largest |k Z_w|, at the outer edge (L5, L7)
\newcommand{\dwAmpPercent}{3}   %median |xi - L_L/x|/|Z_w| - 1 over the steady stretch, percent
\newcommand{\dwKcontPercent}{3.6}   %median k/k_continuum - 1 over the steady stretch, percent (L5)
\newcommand{\dwKlatticePercent}{5.0}   %median k_lattice/k_continuum - 1 over the steady stretch, percent (lattice_wavenumber)
\newcommand{\dwSawtooth}{4\times10^{-7}}   %largest sawtooth at the run's end over |Z_w|, rounded up
\newcommand{\dwNparticles}{32}   %inputs.py particles_per_streamline
\newcommand{\dwForcedMedianPercent}{1}   %median |xi - L7|/|L7| inside the resonance zone (x < -2 eps^1/2), percent
\newcommand{\dwForcedPercent}{3}   %largest |xi - L7|/|L7| inside the resonance zone, percent
\newcommand{\dwTorqueRatio}{0.993}   %mean T/T_L over the last half of the run (L8)
\newcommand{\dwTorqueScatter}{22}   %largest |T/T_L - 1| over the last half, percent, rounded up
\newcommand{\dwXvisALow}{0.0106}   %run A, 400 streamlines: x_vis of the n = 3 fit (measured_damping)
\newcommand{\dwXvisCLow}{0.00689}   %run C, 400 streamlines: x_vis of the n = 3 fit (measured_damping)
\newcommand{\dwXvisDLow}{0.00637}   %run D, 400 streamlines: x_vis of the n = 3 fit (measured_damping)
\newcommand{\dwXvisA}{0.0110}   %run A, 800 streamlines: x_vis of the n = 3 fit (measured_damping)
\newcommand{\dwXvisTheoryA}{0.0110}   %run A: O6 x_vis, beta = 0
\newcommand{\dwXvisShuA}{0.00497}   %run A: Shu (1984) x_vis, beta = -1
\newcommand{\dwNuEffA}{0.304}   %run A: nu_eff/nu (O6)
\newcommand{\dwXvisC}{0.00715}   %run C, 800 streamlines: x_vis of the n = 3 fit (measured_damping)
\newcommand{\dwXvisTheoryC}{0.00743}   %run C: O6 x_vis, beta = 0
\newcommand{\dwXvisShuC}{0.00743}   %run C: Shu (1984) x_vis, beta = -1
\newcommand{\dwNuEffC}{1.00}   %run C: nu_eff/nu (O6)
\newcommand{\dwXvisD}{0.00660}   %run D, 800 streamlines: x_vis of the n = 3 fit (measured_damping)
\newcommand{\dwXvisTheoryD}{0.00680}   %run D: O6 x_vis, beta = 0
\newcommand{\dwXvisShuD}{0.00455}   %run D: Shu (1984) x_vis, beta = -1
\newcommand{\dwNuEffD}{1.30}   %run D: nu_eff/nu (O6)
\newcommand{\dwXvisCHigh}{0.00730}   %run C, 1600 streamlines: x_vis of the n = 3 fit (measured_damping)
\newcommand{\dwNuVisc}{2.7\times10^{-12}}   %density_analysis NU_VISCOUS
\newcommand{\dwShearCoef}{-0.696}   %10/3 - 4 (Omega/kappa)^2 at r_L (O6, beta = 0)
\newcommand{\dwExponentMin}{2.9}   %smallest free-exponent fit n over the viscous runs
\newcommand{\dwExponentMax}{3.1}   %largest free-exponent fit n over the viscous runs
\newcommand{\dwDampFitLow}{10}   %DAMPING_FIT lower ratio, percent
\newcommand{\dwDampFitHigh}{90}   %DAMPING_FIT upper ratio, percent
\newcommand{\dwXvisExtA}{0.0115}   %run A: 2 x_vis(800) - x_vis(400)
\newcommand{\dwXvisExtPercentA}{4}   %run A: |extrapolated/theory - 1|, percent
\newcommand{\dwXvisExtC}{0.00742}   %run C: 2 x_vis(800) - x_vis(400)
\newcommand{\dwXvisExtPercentC}{0.08}   %run C: |extrapolated/theory - 1|, percent
\newcommand{\dwXvisExtD}{0.00684}   %run D: 2 x_vis(800) - x_vis(400)
\newcommand{\dwXvisExtPercentD}{0.6}   %run D: |extrapolated/theory - 1|, percent
\newcommand{\dwConvergenceRatio}{1.8}   %run C: (x_400 - x_800)/(x_800 - x_1600)
